\documentclass[10pt,a4paper]{amsart}
\usepackage{amsmath,amssymb,amsthm,graphicx}
% For commutative diagrams we use Xy-pic
\usepackage[all]{xy}

% We use 2cell for 2-commutative diagrams.
\xyoption{2cell}
% Theorem environments.
%
\theoremstyle{plain}
\newtheorem{theorem}[subsection]{Theorem}
\newtheorem{proposition}[subsection]{Proposition}
\newtheorem{lemma}[subsection]{Lemma}

\theoremstyle{definition}
\newtheorem{definition}[subsection]{Definition}
\newtheorem{example}[subsection]{Example}
\newtheorem{exercise}[subsection]{Exercise}
\newtheorem{situation}[subsection]{Situation}

\theoremstyle{remark}
\newtheorem{remark}[subsection]{Remark}
\newtheorem{remarks}[subsection]{Remarks}

\numberwithin{equation}{subsection}

% Macros
%
\def\lim{\mathop{\mathrm{lim}}\nolimits}
\def\colim{\mathop{\mathrm{colim}}\nolimits}
\def\Spec{\mathop{\mathrm{Spec}}}
\def\Hom{\mathop{\mathrm{Hom}}\nolimits}
\def\Ext{\mathop{\mathrm{Ext}}\nolimits}
\def\SheafHom{\mathop{\mathcal{H}\!\mathit{om}}\nolimits}
\def\SheafExt{\mathop{\mathcal{E}\!\mathit{xt}}\nolimits}
\def\Sch{\mathit{Sch}}
\def\Mor{\mathop{\mathrm{Mor}}\nolimits}
\def\Ob{\mathop{\mathrm{Ob}}\nolimits}
\def\Sh{\mathop{\mathit{Sh}}\nolimits}
\def\NL{\mathop{N\!L}\nolimits}
\def\CH{\mathop{\mathrm{CH}}\nolimits}
\def\proetale{{pro\text{-}\acute{e}tale}}
\def\etale{{\acute{e}tale}}
\def\QCoh{\mathit{QCoh}}
\def\Ker{\mathop{\mathrm{Ker}}}
\def\Im{\mathop{\mathrm{Im}}}
\def\Coker{\mathop{\mathrm{Coker}}}
\def\Coim{\mathop{\mathrm{Coim}}}

% Boxtimes
%
\DeclareMathSymbol{\boxtimes}{\mathbin}{AMSa}{"02}

%
% Macros for moduli stacks/spaces
%
\def\QCohstack{\mathcal{QC}\!\mathit{oh}}
\def\Cohstack{\mathcal{C}\!\mathit{oh}}
\def\Spacesstack{\mathcal{S}\!\mathit{paces}}
\def\Quotfunctor{\mathrm{Quot}}
\def\Hilbfunctor{\mathrm{Hilb}}
\def\Curvesstack{\mathcal{C}\!\mathit{urves}}
\def\Polarizedstack{\mathcal{P}\!\mathit{olarized}}
\def\Complexesstack{\mathcal{C}\!\mathit{omplexes}}
% \Pic is the operator that assigns to X its picard group, usage \Pic(X)
% \Picardstack_{X/B} denotes the Picard stack of X over B
% \Picardfunctor_{X/B} denotes the Picard functor of X over B
\def\Pic{\mathop{\mathrm{Pic}}\nolimits}
\def\Picardstack{\mathcal{P}\!\mathit{ic}}
\def\Picardfunctor{\mathrm{Pic}}
\def\Deformationcategory{\mathcal{D}\!\mathit{ef}}
\usepackage{cleveref}
\crefname{originaldefinition}{original definition}{original definitions}
\Crefname{originaldefinition}{Original definition}{Original definitions}
\makeatletter
\newlabel{situation-setup}{{02QL}{source}{}{sourcepointer-1}{}}
\newlabel{situation-setup@cref}{{[originaldefinition][0][]02QL}{[1][0][]source}{}{}{}}
\makeatother
\crefname{originallemma}{original lemma}{original lemmas}
\Crefname{originallemma}{Original lemma}{Original lemmas}
\makeatletter
\newlabel{section-pre-derived}{{0ESY}{source}{}{sourcepointer-2}{}}
\newlabel{section-pre-derived@cref}{{[originallemma][0][]0ESY}{[1][0][]source}{}{}{}}
\makeatother
\crefname{originallemma}{original lemma}{original lemmas}
\Crefname{originallemma}{Original lemma}{Original lemmas}
\makeatletter
\newlabel{lemma-gysin-at-infty}{{0F9J}{source}{}{sourcepointer-3}{}}
\newlabel{lemma-gysin-at-infty@cref}{{[originallemma][0][]0F9J}{[1][0][]source}{}{}{}}
\makeatother
\crefname{originallemma}{original lemma}{original lemmas}
\Crefname{originallemma}{Original lemma}{Original lemmas}
\makeatletter
\newlabel{lemma-silly}{{0F9H}{source}{}{sourcepointer-4}{}}
\newlabel{lemma-silly@cref}{{[originallemma][0][]0F9H}{[1][0][]source}{}{}{}}
\makeatother
\crefname{originallemma}{original lemma}{original lemmas}
\Crefname{originallemma}{Original lemma}{Original lemmas}
\makeatletter
\newlabel{lemma-push-proper-bivariant}{{0EPK}{source}{}{sourcepointer-5}{}}
\newlabel{lemma-push-proper-bivariant@cref}{{[originallemma][0][]0EPK}{[1][0][]source}{}{}{}}
\makeatother
\crefname{originallemma}{original lemma}{original lemmas}
\Crefname{originallemma}{Original lemma}{Original lemmas}
\makeatletter
\newlabel{lemma-silly-silly}{{0F9I}{source}{}{sourcepointer-6}{}}
\newlabel{lemma-silly-silly@cref}{{[originallemma][0][]0F9I}{[1][0][]source}{}{}{}}
\makeatother
\title{Existence of the localized Chern class as a bivariant class on a closed subscheme, with the pushforward identity onto the zero fibre}
\author{The Stacks Project Authors}
\date{stacks-project:chow.tex:local-snapshot-20261012; GNU FDL 1.2 or later}
\begin{document}
\maketitle
\noindent\textit{Editorial note.} All delivered mathematical text is the author's own, copied from the pinned Stacks Project chapter file chow.tex; no statement or proof step was rewritten, reordered or removed.\par\medskip
\noindent\textit{Editorial note.} Preamble provenance: the pinned chapter file opens with an input of the shared preamble and carries no preamble text of its own, so the project's single Stacks-wide preamble staged at sources/stacks-chow/preamble.tex is excerpted verbatim as two spans, lines 14-18 (Xy-pic setup) and lines 154-217 (the native theorem declarations and mathematical macros). Nothing was added, removed or redefined; the wrapper class and the amsmath/amssymb/amsthm/graphicx base packages are supplied by the assembler.\par\medskip
\noindent\textit{Editorial note.} The counted statement and proof contain no cross reference, no citation and no included file, so the delivered body is self contained.\par\medskip
\noindent\textit{Editorial note.} Canonical location of the counted proposition: Stacks tag 0F9K (sources/stacks-chow/pinned-original-source.tex line 10204).\par\medskip
\noindent\textit{Allowed standard prerequisites (editorial).}\begin{itemize}
\item The Gysin-at-infinity bivariant class C and its defining property. Same chapter, Lemma labelled lemma-gysin-at-infty (source line 10026), invoked at proof lines 10236 and 10267. Provides C in A\textasciicircum{}0(W\_infinity -> X) with i\_infinity,*(C cap alpha) = i\_\{0,*\}alpha, which is the ingredient that moves the construction from W\_infinity down to X.
\item The local construction of the class c'\_p(Q|\_E) on a closed subscheme where the module is zero or of small rank. Same chapter, Lemma labelled lemma-silly (source line 9695), invoked at proof lines 10237-10239. Supplies c'\_p(Q|\_E) in A\textasciicircum{}p(E -> W\_infinity), applicable because of assumption (A2); the assembly of the three classes is inside the counted span.
\item Proper pushforward of bivariant classes and the comparison lemma for the restricted class. Same chapter, Lemma labelled lemma-push-proper-bivariant (source line 6167) at proof line 10247 and Lemma labelled lemma-silly-silly (source line 9792) at proof line 10265. Supply that the composite defining c'\_p(Q) is again a bivariant class and the comparison with c\_p(Q|\_\{W\_infinity\}).
\item Chern classes of a perfect object and the base situation (S, delta). Same chapter, Section labelled section-pre-derived (source line 9042) and Situation labelled situation-setup (source line 1372), invoked in the statement. Fix assumption (A0) and the ambient operational Chow groups.
\end{itemize}\medskip
\ifdefined\hypertarget\hypertarget{sourcepointer-1}{}\fi
\noindent\textit{Source reference.} Original definition 02QL: sources/stacks-chow/pinned-original-source.tex, source lines 1372--1372. Same-chapter pointer into the author environment situation at source line 1372; invoked at proof line None of the counted proof. It is retained only as a navigation pointer; every load-bearing step is inside the delivered span.\par\medskip
\ifdefined\hypertarget\hypertarget{sourcepointer-2}{}\fi
\noindent\textit{Source reference.} Original lemma 0ESY: sources/stacks-chow/pinned-original-source.tex, source lines 9042--9042. Same-chapter pointer into the author environment lemma at source line 9042; invoked at proof line None of the counted proof. It is retained only as a navigation pointer; every load-bearing step is inside the delivered span.\par\medskip
\ifdefined\hypertarget\hypertarget{sourcepointer-3}{}\fi
\noindent\textit{Source reference.} Original lemma 0F9J: sources/stacks-chow/pinned-original-source.tex, source lines 10026--10026. Same-chapter pointer into the author environment lemma at source line 10026; invoked at proof line 10237 of the counted proof. It is retained only as a navigation pointer; every load-bearing step is inside the delivered span.\par\medskip
\ifdefined\hypertarget\hypertarget{sourcepointer-4}{}\fi
\noindent\textit{Source reference.} Original lemma 0F9H: sources/stacks-chow/pinned-original-source.tex, source lines 9695--9695. Same-chapter pointer into the author environment lemma at source line 9695; invoked at proof line 10239 of the counted proof. It is retained only as a navigation pointer; every load-bearing step is inside the delivered span.\par\medskip
\ifdefined\hypertarget\hypertarget{sourcepointer-5}{}\fi
\noindent\textit{Source reference.} Original lemma 0EPK: sources/stacks-chow/pinned-original-source.tex, source lines 6167--6167. Same-chapter pointer into the author environment lemma at source line 6167; invoked at proof line 10247 of the counted proof. It is retained only as a navigation pointer; every load-bearing step is inside the delivered span.\par\medskip
\ifdefined\hypertarget\hypertarget{sourcepointer-6}{}\fi
\noindent\textit{Source reference.} Original lemma 0F9I: sources/stacks-chow/pinned-original-source.tex, source lines 9792--9792. Same-chapter pointer into the author environment lemma at source line 9792; invoked at proof line 10265 of the counted proof. It is retained only as a navigation pointer; every load-bearing step is inside the delivered span.\par\medskip
\begin{lemma}
\label{lemma-localized-chern-pre}
Let $(S, \delta)$ be as in Situation \ref{situation-setup}. Let $X$ be
locally of finite type over $S$. Let $Z \subset X$ be a closed subscheme.
Let
$$
b : W \longrightarrow \mathbf{P}^1_X
$$
be a proper morphism of schemes. Let $Q \in D(\mathcal{O}_W)$ be a
perfect object. Denote $W_\infty \subset W$ the inverse image of the divisor
$D_\infty \subset \mathbf{P}^1_X$ with complement $\mathbf{A}^1_X$.
We assume
\begin{enumerate}
\item[(A0)] Chern classes of $Q$ are defined
(Section \ref{section-pre-derived}),
\item[(A1)] $b$ is an isomorphism over $\mathbf{A}^1_X$,
\item[(A2)] there exists a closed subscheme $T \subset W_\infty$
containing all points of $W_\infty$ lying over $X \setminus Z$ such that
$Q|_T$ is zero, resp.\ isomorphic to a finite locally free
$\mathcal{O}_T$-module of rank $< p$ sitting in cohomological degree $0$.
\end{enumerate}
Then there exists a canonical bivariant class
$$
P'_p(Q),\text{ resp. }c'_p(Q) \in A^p(Z \to X)
$$
with
$(Z \to X)_* \circ P'_p(Q) = P_p(Q|_{X \times \{0\}})$,
resp.\ $(Z \to X)_* \circ c'_p(Q) = c_p(Q|_{X \times \{0\}})$.
\end{lemma}

\begin{proof}
Denote $E \subset W_\infty$ the inverse image of $Z$. Then
$W_\infty = T \cup E$ and $b$ induces a proper morphism $E \to Z$.
Denote $C \in A^0(W_\infty \to X)$ the bivariant class constructed
in Lemma \ref{lemma-gysin-at-infty}. Denote $P'_p(Q|_E)$, resp.\ $c'_p(Q|_E)$
in $A^p(E \to W_\infty)$ the bivariant class constructed
in Lemma \ref{lemma-silly}. This makes sense because
$(Q|_E)|_{T \cap E}$ is zero, resp.\ isomorphic to a finite locally free
$\mathcal{O}_{E \cap T}$-module of rank $< p$ sitting in
cohomological degree $0$ by assumption (A2). Then we define
$$
P'_p(Q) = (E \to Z)_* \circ P'_p(Q|_E) \circ C,\text{ resp. }
c'_p(Q) = (E \to Z)_* \circ c'_p(Q|_E) \circ C
$$
This is a bivariant class, see Lemma \ref{lemma-push-proper-bivariant}.
Since $E \to Z \to X$ is equal to $E \to W_\infty \to W \to X$ we see that
\begin{align*}
(Z \to X)_* \circ c'_p(Q)
& =
(W \to X)_* \circ i_{\infty, *} \circ (E \to W_\infty)_*
\circ c'_p(Q|_E) \circ C \\
& =
(W \to X)_* \circ i_{\infty, *} \circ c_p(Q|_{W_\infty}) \circ C \\
& =
(W \to X)_* \circ c_p(Q) \circ i_{\infty, *} \circ C \\
& =
(W \to X)_*\circ c_p(Q) \circ i_{0, *} \\
& =
(W \to X)_* \circ i_{0, *} \circ c_p(Q|_{X \times \{0\}}) \\
& =
c_p(Q|_{X \times \{0\}})
\end{align*}
The second equality holds by Lemma \ref{lemma-silly-silly}.
The third equality because $c_p(Q)$ is a bivariant class.
The fourth equality by Lemma \ref{lemma-gysin-at-infty}.
The fifth equality because $c_p(Q)$ is a bivariant class.
The final equality because $(W_0 \to W) \circ (W \to X)$
is the identity on $X$ if we identify $W_0$ with $X$ as we've
done above. The exact same sequence of equations works to
prove the property for $P'_p(Q)$.
\end{proof}

\end{document}
